\documentclass[11pt]{article}
\usepackage[utf8]{inputenc}
\usepackage{amsmath,amssymb}
\usepackage{graphicx}
\usepackage{hyperref}
\usepackage{booktabs}
\usepackage[margin=1in]{geometry}
\usepackage{xcolor}
\usepackage{enumitem}

\definecolor{huimai-blue}{RGB}{0, 102, 204}

\title{A General-Purpose Agent Collaboration Framework \\
for Domain-Specific AI Systems: \\
From Shanzhai Methodology to the Doctor's Framework}

\author{SharpAgent Research \\
\texttt{research@bio-shandonghuiumai.com}}

\date{June 2026}

\begin{document}

\maketitle

\begin{abstract}
We present a general-purpose agent collaboration framework designed for 
constructing domain-specific AI systems, with a validated implementation 
in clinical decision support. The core contribution is a three-layer 
agent architecture---strategist, architect, and executor---that formalizes 
a problem-driven development methodology we term ``Shanzhai with Structure.'' 
We demonstrate the framework's effectiveness through the Doctor's Framework, 
a production clinical knowledge system featuring a 68-module knowledge base, 
a jurisdiction-aware tripartite labeling system (Z0 layer), and a structured 
guideline index layer (GuideLayer) covering 102 clinical references across 
26 medical specialties. The system has been deployed in collaboration with 
a tertiary hospital in China, validating the framework's transition from 
abstract collaboration model to operational AI infrastructure. We argue 
that structured agent collaboration, rather than monolithic model training, 
is the scalable path to domain-specific AI that is auditable, updatable, 
and jurisdiction-aware.
\end{abstract}

\section{Introduction}

Large language models (LLMs) have demonstrated remarkable capabilities in 
general-domain question answering \cite{brown2020language}, but their deployment 
in high-stakes professional domains remains problematic. Three structural 
limitations persist: (1) LLMs lack verifiable, updatable domain knowledge 
bases, operating instead on frozen training corpora; (2) they cannot reliably 
distinguish jurisdiction-specific standards (e.g., Chinese CSCO guidelines 
vs. American NCCN guidelines for the same cancer); and (3) they provide no 
audit trail for clinical decisions, making them legally and ethically 
untenable in regulated environments.

The prevailing response to these limitations has been to invest in 
application-specific model fine-tuning, retrieval-augmented generation (RAG) 
\cite{lewis2020retrieval}, or prompt engineering. These approaches, while 
valuable, share a common weakness: they embed domain knowledge into model 
parameters or narrow retrieval pipelines, making systematic updates, 
jurisdiction switching, and cross-guideline conflict detection fundamentally 
difficult.

We propose an alternative paradigm: a \textbf{general-purpose agent 
collaboration framework} that separates domain knowledge from reasoning, 
and organizes AI system construction around a structured three-layer 
agent model. This framework is not an application but a \textit{methodology 
for building applications}---a design pattern for domain-specific AI that 
prioritizes auditability, updatability, and jurisdiction-awareness over 
raw generative capability.

\subsection{Contributions}

This paper makes three contributions:

\begin{enumerate}[leftmargin=*]
    \item \textbf{Formalization of the Sharp Framework}: a three-layer agent 
    collaboration model (strategist, architect, executor) that provides a 
    reusable architecture for domain-specific AI construction.
    \item \textbf{The Doctor's Framework}: a complete implementation of the 
    Sharp Framework in clinical medicine, featuring a 68-module knowledge 
    base organized in four scientific layers (L1--L4), a jurisdiction-aware 
    Z0 tripartite labeling layer, and a structured guideline index layer 
    (GuideLayer) with 102 clinical references across 26 specialties.
    \item \textbf{Validation}: deployment in collaboration with a tertiary 
    hospital in Shandong Province, China, demonstrating the framework's 
    transition from abstract collaboration model to operational clinical 
    infrastructure.
\end{enumerate}

\section{The Sharp Framework}

\subsection{Philosophy: ``Shanzhai with Structure''}

The term \textit{Shanzhai} (山寨) carries a specific cultural meaning in 
Chinese innovation discourse: pragmatic, resource-constrained problem-solving 
that prioritizes working solutions over theoretical elegance. When 
structured---paired with clear role boundaries, verification protocols, 
and systematic knowledge management---this approach becomes a powerful 
methodology for AI system construction.

The Sharp Framework formalizes this intuition into three design principles:

\begin{enumerate}[leftmargin=*]
    \item \textbf{Problem-driven development}: Build only what the problem 
    demands. Defer abstraction until pattern recognition is justified by 
    multiple concrete instances.
    \item \textbf{Role separation with verification}: Separate strategic 
    decision-making, architectural coordination, and technical execution 
    into distinct agent roles with explicit verification at each boundary.
    \item \textbf{Knowledge externalization}: All domain knowledge lives 
    in auditable, version-controlled external structures (knowledge bases, 
    guideline indices, jurisdiction tables) rather than in model parameters.
\end{enumerate}

\subsection{Three-Layer Agent Architecture}

\begin{figure}[h]
\centering
\begin{verbatim}
Strategic Layer (Human-in-the-loop)
    ↓  Vision, trust allocation, go/no-go decisions
Architect Layer (cn001 — coordination agent)
    ↓  Architecture design, task decomposition, quality control
Executor Layer (hk001, cn002, ... — specialist agents)
    ↓  Implementation, testing, cross-review, deployment
\end{verbatim}
\caption{The Sharp Framework three-layer agent architecture.}
\label{fig:three-layer}
\end{figure}

The three layers operate with clearly defined scopes and verification 
protocols:

\textbf{Strategic Layer.} A human decision-maker provides vision, domain 
expertise, and final authority on trust-critical decisions. The strategist 
does not write code or review architectures---only allocates trust and 
makes go/no-go judgments. This is not a ceremonial role; it is the 
locus of ``what should we build'' and ``is this safe to deploy.''

\textbf{Architect Layer.} A coordination agent (cn001) translates strategic 
intent into technical specifications, decomposes tasks, verifies 
cross-layer consistency, and maintains the global knowledge architecture. 
This agent does not execute implementation tasks---only designs, coordinates, 
and verifies. All communication between layers passes through this agent 
via a formal inbox protocol.

\textbf{Executor Layer.} Specialist agents (hk001 for cloud operations, 
cn002 for knowledge base engineering, and others) receive well-defined 
task specifications from the architect layer, execute implementation, 
perform cross-review, and report results back through the formal 
communication channel. Each executor operates within a bounded scope 
defined by its task contract.

\subsection{Formal Communication Protocol}

To prevent information leakage and role confusion, the Sharp Framework 
mandates a formal communication protocol:

\begin{itemize}[leftmargin=*]
    \item All inter-agent communication passes through designated inbox 
    directories (local filesystem for co-located agents, SCP for remote agents).
    \item Each message follows a standardized format: sender, recipient, 
    task ID, payload, verification criteria.
    \item No agent may access another agent's filesystem or working directory.
    \item After message delivery, the sender notifies the recipient via an 
    agreed side channel.
\end{itemize}

This protocol is not an optimization---it is a safety boundary. By making 
all communication explicit, auditable, and replayable, we eliminate the 
class of errors caused by implicit assumptions about what another agent 
``knows'' or ``has done.''

\section{The Doctor's Framework: A Validated Implementation}

The Doctor's Framework is a complete instantiation of the Sharp Framework 
for clinical decision support. It is not a model---it is a system 
architecture that orchestrates knowledge retrieval, jurisdiction-aware 
reasoning, and guideline-aware recommendation.

\subsection{Knowledge Architecture}

\begin{figure}[h]
\centering
\begin{verbatim}
┌─────────────────────────────────────────────┐
│              Knowledge Base (68 modules)     │
│  L1: Basic Sciences (8)    Anatomy, Biochem  │
│  L2: Core Applications (13)  Pathophys, Pharm│
│  L3: Clinical Cover (21)   Emergency, Surg   │
│  L4: Prospective (7)       AI Dx, Genomics   │
│  T:  Traditional Medicine (13)  TCM, Ayurveda│
│  E:  Environmental Exposure (6)  NBC Warfare │
├─────────────────────────────────────────────┤
│         Z0: Jurisdiction Layer               │
│  Tripartite Label: [T_bio, T_timeline, T_medhist] │
│  Migration Baseline Function                 │
├─────────────────────────────────────────────┤
│         GuideLayer: Guideline Index          │
│  102 clinical references, 26 specialties     │
│  Jurisdiction-tagged, cross-referenced       │
│  Physically isolated from Knowledge Base     │
└─────────────────────────────────────────────┘
\end{verbatim}
\caption{The Doctor's Framework knowledge architecture.}
\label{fig:knowledge}
\end{figure}

The knowledge architecture is organized in four distinct layers, each 
with a clearly defined scope:

\textbf{L1--L4: Scientific Knowledge Base.} The core scientific knowledge 
is organized in four layers of increasing clinical specificity. L1 covers 
basic sciences (biochemistry, genomics, immunology); L2 covers core 
clinical applications (pathophysiology, pharmacology); L3 covers clinical 
specialties (emergency medicine, surgery, pediatrics); L4 covers prospective 
domains (AI-assisted diagnosis, pharmacogenomics). Additionally, T covers 
traditional medicine systems from 12 cultural traditions, and E covers 
environmental and occupational exposure.

Critically, the knowledge base contains \textbf{no clinical guidelines}. 
It stores scientific facts---anatomy, physiology, pathology, pharmacology, 
epidemiology---that are globally valid and jurisdiction-independent.

\textbf{Z0: Jurisdiction Layer.} Before any clinical query reaches the 
knowledge base, it passes through Z0, which constructs a tripartite label:
\begin{equation}
\text{Label} = [T_{\text{bio}}, T_{\text{timeline}}, T_{\text{medhist}}]
\end{equation}
where $T_{\text{bio}}$ encodes the patient's biological jurisdiction 
(genetic ancestry, birthplace), $T_{\text{timeline}}$ encodes migration 
history and residence duration, and $T_{\text{medhist}}$ encodes the 
jurisdictional origin of past medical interventions.

Z0 also implements a \textbf{migration baseline function} that models how 
physiological parameters drift when a patient moves between jurisdictions. 
For example, a patient of Chinese ancestry who has lived in the United 
States for 15 years will have BMI thresholds, salt intake baselines, and 
vitamin D reference ranges that are partially Chinese, partially American, 
and time-dependent.

\textbf{GuideLayer.} Clinical guidelines (CSCO, NCCN, ESC, expert consensus 
documents) are stored in a physically separate index, tagged by jurisdiction, 
and cross-referenced for conflicts. The GuideLayer contains 102 structured 
records across 26 specialties at the time of writing. When the system 
generates a recommendation, it explicitly cites which guideline(s) informed 
it, what jurisdiction they apply to, and whether conflicting guidance 
exists in other jurisdictions.

The physical separation of GuideLayer from the knowledge base is a 
deliberate architectural decision. Guidelines are \textit{interpretations} 
of science within a specific regulatory and cultural context---not science 
itself. Mixing them would pollute the globally valid knowledge base with 
jurisdiction-specific opinions.

\subsection{Clinical Workflow}

A typical clinical query follows this path:

\begin{enumerate}[leftmargin=*]
    \item Patient data enters through a standardized questionnaire 
    (jurisdiction-relevant questions only).
    \item Z0 constructs the tripartite label and computes the mixed 
    jurisdiction weight vector.
    \item The knowledge index retrieves relevant modules by symptom and 
    disease matching.
    \item The GuideLayer retrieves applicable guidelines filtered by the 
    patient's jurisdiction vector.
    \item The system generates a recommendation with explicit citation 
    of (a) which scientific modules informed the reasoning, (b) which 
    guidelines were applied, and (c) what jurisdiction they represent.
\end{enumerate}

The output is not a black-box prediction---it is a structured, auditable 
recommendation that a human physician can verify against the cited sources.

\subsection{GuideLayer: A Structured Guideline Index}

The GuideLayer represents a methodological innovation in clinical AI: a 
jurisdiction-tagged, cross-referenced index of clinical practice guidelines 
maintained separately from scientific knowledge.

At the time of writing, the GuideLayer contains 102 records, including:

\begin{itemize}[leftmargin=*]
    \item 7 CSCO (Chinese Society of Clinical Oncology) guidelines for 
    major cancer types, with full structured extraction of staging systems, 
    first-line and second-line regimens, and cross-references to NCCN/ESMO 
    guidelines.
    \item 4 Chinese Medical Association clinical guidelines and expert 
    consensus documents.
    \item 87 additional Chinese expert consensus documents indexed by 
    title, society, year, and specialty.
    \item 4 hospital formulary records from tertiary hospitals in Shandong 
    Province.
\end{itemize}

Each guideline record includes: a unique jurisdiction-tagged ID, the 
issuing society, disease scope, year, evidence grading system, structured 
treatment recommendations (regimen, biomarker, evidence level, clinical 
trial reference), and cross-references marking conflicts and alignments 
with guidelines from other jurisdictions.

The GuideLayer is designed for incremental maintenance: as guidelines are 
updated annually, new records supersede old ones while preserving the 
update history, creating a longitudinal record of clinical consensus 
evolution.

\subsection{Technical Implementation}

The system is implemented as a Python/FastAPI backend with a pure HTML 
frontend (no framework dependency), deployed on Alibaba Cloud ECS behind 
Cloudflare Tunnel. The knowledge index uses substring sliding-window 
matching augmented with a disease-specific terminology dictionary. 

The software architecture follows the Sharp Framework's role separation:
\begin{itemize}[leftmargin=*]
    \item \textbf{Strategist (Lao Wang)}: clinical domain validation, 
    hospital partnership, go/no-go decisions on deployment.
    \item \textbf{Architect (cn001)}: knowledge architecture design, 
    Guideline index maintenance, quality verification of extractions.
    \item \textbf{Executor (cn002)}: knowledge base module authoring, 
    benchmark testing, OCR pipeline for guideline digitization.
    \item \textbf{Executor (hk001)}: cloud deployment, API stability, 
    cross-review of extracted content.
\end{itemize}

\subsection{Benchmark Results}

The knowledge retrieval system was benchmarked on a set of clinical 
queries spanning all 68 knowledge modules. Initial results showed a 
retrieval hit rate of 13\%, which after three rounds of pipeline debugging 
improved to 76.7\% (content-level hit rate). The remaining 23.3\% gap is 
concentrated in 10 modules (4 respiratory/emergency, 6 others) that require 
additional content authoring.

The three root causes identified and fixed during benchmarking illustrate 
the framework's strength: (1) a Unicode ghost character in the API key 
causing silent degradation; (2) a path break in the local content pipeline 
where distilled module content was read but never passed to the LLM; and 
(3) a Chinese text matching algorithm that required substring sliding-window 
replacement. Each failure mode was traced to a specific architectural 
component, fixed, and verified---an audit trail impossible in monolithic 
fine-tuned models.

\subsection{Structured Knowledge Extraction Pipeline}

A production-grade OCR pipeline was developed using Tencent Cloud's 
GeneralBasicOCR API (TC3-HMAC-SHA256 signing). The pipeline processed 
181 source files (143 DOCX with 1,206 embedded images, 7 PDFs totaling 
1,317 pages, 9 WeChat screenshots, and 10 Excel/PPT files) across three
batches, producing 998 OCR calls with a 97\% success rate, yielding
2,548,459 characters of structured clinical text spanning CSCO oncology
guidelines, CMA clinical practice guidelines, and pharmaceutical formulary
data. An additional 79 WeChat-sourced pharmaceutical screenshots were
processed to recover drug catalogue data blocked by WeChat anti-scraping,
demonstrating the framework's resilience to real-world data acquisition
barriers.

\subsection{Pharmaceutical Knowledge Base}

A structured drug knowledge base was constructed covering 3,369 drug
records from 11 rounds of Chinese national centralized procurement
(spanning batch 1 through 10, plus the insulin special procurement),
merged with hospital drug formularies from Tai'an Central Hospital,
Jining First People's Hospital, and Linyi People's Hospital.
Each record includes generic name, dosage form, specifications,
manufacturer, procurement batch, and pricing data.

\subsection{Standard Alignment Audit}

A comprehensive gap audit was conducted against two recently published
Chinese national consensus documents on AI-enabled laboratory report
interpretation: (1) the Expert Consensus on Intelligent Agent-Based
Laboratory Report Interpretation (2026), and (2) the Expert Consensus on
Standard Dataset Elements for Medical Laboratory Science (2026). The
audit confirmed that the Doctor's Framework's Five-Element Audit System
(Reviewer-Auditor Verifier Calibration Chain) exceeds the consensus
requirement for human-in-the-loop design through its independent reviewer
bias coefficient. Two critical gaps were identified---LOINC/UCUM/SNOMED
CT/FHIR international standard mapping and data privacy governance---now
scheduled for v0.4 implementation.

\section{Related Work}

\textbf{Agent Frameworks.} AutoGen \cite{wu2023autogen}, CrewAI, and 
LangGraph provide general-purpose multi-agent orchestration. The Sharp 
Framework differs in emphasizing role separation with explicit 
verification boundaries and a formal communication protocol, making it 
more suitable for regulated domains.

\textbf{Clinical AI.} Med-PaLM 2 \cite{singhal2023towards} and GPT-4-based 
clinical systems demonstrate impressive generative capabilities but suffer 
from the jurisdiction-ambiguity problem: they cannot reliably distinguish 
CSCO from NCCN guidelines. Retrieval-augmented approaches \cite{zhu2024rag} 
partially address this but typically embed retrieval into model pipelines 
rather than maintaining a separate, auditable knowledge architecture.

\textbf{Knowledge-Grounded Systems.} The separation of knowledge from 
reasoning has been advocated in symbolic AI \cite{lenat1990cyc} and 
more recently in neuro-symbolic systems \cite{garcez2020neurosymbolic}. 
The Doctor's Framework contributes a practical instantiation of this 
principle in a deployed clinical system with a concrete knowledge 
architecture and a validated retrieval pipeline.

\section{Limitations and Future Work}

\textbf{Knowledge Coverage.} The 68-module knowledge base, while 
comprehensive in structure, has uneven content depth across modules. 
Ten modules remain below the 70\% hit rate threshold and require 
additional content authoring by domain specialists.

\textbf{Guideline Completeness.} The GuideLayer currently covers Chinese
(CN) clinical guidelines (181 source files processed via OCR, 102
indexed guideline records in index.json with 11 deep-extracted CSCO
oncology + CMA guideline full texts). US (NCCN), European (ESC/ESMO),
Japanese, and Southeast Asian guidelines are planned but not yet indexed.
Concurrently, a pharmaceutical knowledge base of 3,369 drug records from
11 rounds of national centralized procurement has been constructed, with
top-200 package insert extraction (indications, dosage, contraindications)
underway for clinical decision support augmentation.

\textbf{Evaluation.} The current evaluation is limited to retrieval
accuracy benchmarks. Clinical outcome evaluation through the hospital
partnership is planned but has not yet been conducted. An independent
gap audit against two national AI-in-laboratory-medicine consensus
documents has been completed, confirming human-in-the-loop compliance
and identifying LOINC/UCUM/SNOMED CT mapping as the top priority for
hospital integration.

\textbf{Scaling.} The framework's reliance on human strategist input 
and formal communication protocols limits the number of concurrent 
projects. Formal verification of the communication protocol and 
partial automation of the architect layer are directions for future 
work.

\section{Conclusion}

We have presented the Sharp Framework, a general-purpose agent 
collaboration methodology for domain-specific AI construction, and its 
validated implementation in the Doctor's Framework for clinical decision 
support. The key architectural insight is the separation of scientific 
knowledge, jurisdiction-specific parameters, and clinical practice 
guidelines into three physically distinct layers, each with its own 
maintenance cadence and verification protocol. This architecture enables 
a clinical AI system that is auditable (every recommendation cites its 
sources), updatable (guidelines and knowledge modules can be revised 
independently), and jurisdiction-aware (the same disease query produces 
different, appropriately localized recommendations for patients from 
different jurisdictions).

The Sharp Framework is not a product---it is a design pattern. Its value 
lies in its replicability: the same three-layer agent architecture, 
formal communication protocol, and knowledge externalization principles 
can be applied to construct domain-specific AI systems in law, finance, 
engineering, and any other regulated professional domain where 
auditability and jurisdiction-awareness matter.

\section*{Acknowledgments}

We thank the Department of Respiratory Medicine at the partner hospital 
in Shandong Province for clinical domain validation. We acknowledge the 
contributions of cn002 for knowledge base engineering and benchmark 
methodology, and hk001 for cloud infrastructure and cross-review.

\begin{thebibliography}{99}

\bibitem{brown2020language}
T.~Brown et al., ``Language Models are Few-Shot Learners,'' in 
\textit{NeurIPS}, 2020.

\bibitem{lewis2020retrieval}
P.~Lewis et al., ``Retrieval-Augmented Generation for Knowledge-Intensive 
NLP Tasks,'' in \textit{NeurIPS}, 2020.

\bibitem{wu2023autogen}
Q.~Wu et al., ``AutoGen: Enabling Next-Gen LLM Applications via Multi-Agent 
Conversation,'' \textit{arXiv:2308.08155}, 2023.

\bibitem{singhal2023towards}
K.~Singhal et al., ``Towards Expert-Level Medical Question Answering with 
Large Language Models,'' \textit{arXiv:2305.09617}, 2023.

\bibitem{zhu2024rag}
Y.~Zhu et al., ``A Survey of Retrieval-Augmented Generation for Large 
Language Models,'' \textit{arXiv:2312.10997}, 2024.

\bibitem{lenat1990cyc}
D.~B.~Lenat and R.~V.~Guha, \textit{Building Large Knowledge-Based Systems: 
Representation and Inference in the Cyc Project}. Addison-Wesley, 1990.

\bibitem{garcez2020neurosymbolic}
A.~d'Avila Garcez and L.~C.~Lamb, ``Neurosymbolic AI: The 3rd Wave,'' 
\textit{arXiv:2012.05876}, 2020.

\end{thebibliography}

\end{document}
